\section{Conclusion}
\label{sec:conclusion}

This work provides the first systematic evaluation of Generative AI's dual impact on spam detection through 4,080 training-evaluation units spanning two LLM models, three prompt strategies, and two classifier architectures. Our experiments reveal three critical findings: (1) LLM-generated synthetic spam can partially substitute for real training data with dramatic model-dependent quality gaps—GPT-4.1-mini maintains F1 = 0.735-0.827 at 100\% synthetic training while Claude-3.5-Haiku degrades to F1 = 0.457-0.796, establishing model selection as the primary determinant of synthetic data utility; (2) zero-shot detection of AI-generated spam achieves moderate baseline performance (GPT detection: F1 = 0.744-0.831, Claude detection: F1 = 0.532-0.811), with substantial variation by target model, prompt strategy, and classifier architecture; (3) cross-model synthetic augmentation improves detection effectiveness (GPT detection: F1 $\rightarrow$ 0.894, Claude detection: F1 $\rightarrow$ 0.860), with diversity-driven training data providing larger gains than quality-optimized data. Random Forest consistently outperforms SVM across all tracks through superior robustness to distributional shifts, prompt invariance, and cross-model generalization. Our contributions include rigorous empirical foundations for understanding LLM-generated synthetic spam, validation of cross-model augmentation strategies, and identification of architecture-dependent performance patterns. For cybersecurity practitioners, these findings reveal that Generative AI fundamentally favors attackers through low cost and high effectiveness, though defenders can leverage ensemble architectures and cross-model augmentation to maintain detection capabilities. Future work should investigate open-source LLMs, transformer-based classifiers operating on semantic embeddings, and generalization to other cybersecurity domains including network intrusion detection and malware classification.
